\section{最小二乘：为无解方程寻找最佳近似}
\label{sec:03-Linear-Algebra/03-Orthogonality/03}

你手上只有三个点

\[
(1,1),\qquad(2,2),\qquad(3,2),
\]

想用一条直线 \(y=C+Dt\) 穿过它们。动笔列方程，一切都还顺利：

\[
\begin{aligned}
C+D&=1,\\
C+2D&=2,\\
C+3D&=2.
\end{aligned}
\]

前两个方程立刻给出 \(D=1,\;C=0\)。把这对参数送进第三个方程，左边算出来是 \(3\)，右边却是 \(2\)。

这条直线穿过了前两个点，却从第三个点旁边擦过去了。不是算错了——是三个点本来就不共线。手中的模型只有两个参数，却要同时迁就三项包含噪声的观测。

这是个典型的困境：方程组写得出，却没有精确解。

\subsection{从无解转向“残差最小”}

遇到无解方程，一个自然的反应是：能不能退一步，找个“差不多”的解？

把方程写成矩阵：

\[
\underbrace{\begin{pmatrix}
1&1\\
1&2\\
1&3
\end{pmatrix}}_{A}
\underbrace{\begin{pmatrix}C\\  D\end{pmatrix}}_{x} = \underbrace{\begin{pmatrix}1\\2\\2\end{pmatrix}}_{b}.
\]

\(A\) 只有两列，所以无论 \(x\) 取什么值，\(Ax\) 永远落在 \(A\) 的列空间里——\(\R^3\) 中过原点的一个二维平面。而 \(b\) 不在这个平面上。精确解不存在的原因就在这里：观测向量落在模型能够到达的范围之外。

既然精确抵达不可能，那就换个问法：在列空间中，哪个向量离 \(b\) 最近？

对任意参数 \(x\)，定义残差

\[
r=b-Ax.
\]

它的长度衡量预测与观测的差距。最小二乘选择使残差平方长度最小的参数：

\[
\hat x\in\argmin_x\|b-Ax\|_2^2.
\]

取平方而不是直接取长度，纯粹是计算上的便利：平方根不会改变最小值的位置，却省掉了处处携带根号的麻烦，而且让不同观测的误差以平方形式叠加——这正好对应“总误差”的直观。

现在问题从“解方程”变成了“在子空间里找最近点”。上一节已经为这个任务备好了工具：最近点就是正交投影。

\subsection{正交条件直接给出方程}

如果 \(\hat x\) 使 \(A\hat x\) 成为 \(b\) 在列空间上的投影，那么残差

\[
\hat r=b-A\hat x
\]

必须与列空间中的每一个方向垂直。列空间中任意向量都可以写成 \(Ay\)，因此垂直条件就是

\[
\hat r\perp Ay,\quad\forall y.
\]

这等价于 \(\hat r\) 与 \(A\) 的每一列正交，即

\[
A^{\trans}\hat r=0.
\]

展开：

\[
A^{\trans}(b-A\hat x)=0\quad\Longrightarrow\quad A^{\trans}A\hat x=A^{\trans}b.
\]

这就是正规方程。名字听起来有点吓人，含义却很朴素：在最优参数处，残差与所有可能的预测调整方向都正交——任何对参数的微小变动都不能继续缩短残差了。

用勾股关系验证一遍会更踏实。设 \(x=\hat x+h\)，把残差拆成两部分：

\[
b-Ax = b-A\hat x - Ah = \hat r - Ah.
\]

因为 \(\hat r\perp\operatorname{Col}(A)\) 且 \(Ah\in\operatorname{Col}(A)\)，这两项正交。于是

\[
\|b-Ax\|_2^2
= \|\hat r\|_2^2 + \|Ah\|_2^2
\geq \|\hat r\|_2^2.
\]

第二项非负，所以任何偏离 \(\hat x\) 的参数都只会让残差变大。正交性不仅仅是一个几何描述——它精确地刻画了最优性。

\subsection{回到数字：把那条直线算出来}

对本例的数据，

\[
A^{\trans}A=
\begin{pmatrix}
1&1&1\\
1&2&3
\end{pmatrix}
\begin{pmatrix}
1&1\\
1&2\\
1&3
\end{pmatrix}
=
\begin{pmatrix}
3&6\\
6&14
\end{pmatrix},
\qquad
A^{\trans}b=
\begin{pmatrix}
1&1&1\\
1&2&3
\end{pmatrix}
\begin{pmatrix}1\\2\\2\end{pmatrix}
=
\begin{pmatrix}5\\11\end{pmatrix}.
\]

正规方程成为

\[
\begin{aligned}
3\hat C+6\hat D&=5,\\
6\hat C+14\hat D&=11.
\end{aligned}
\]

两倍第一式减第二式：\(-2\hat D=-1\)，得 \(\hat D=\frac12\)；代回得 \(\hat C=\frac23\)。

拟合直线是

\[
\hat y=\frac23+\frac12t.
\]

把参数带回模型，对应三个 \(t=1,2,3\) 的预测值：

\[
\hat b=A\hat x=
\begin{pmatrix}
7/6\\[2pt]5/3\\[2pt]13/6
\end{pmatrix},
\qquad
\hat r=b-\hat b=
\begin{pmatrix}
-1/6\\[2pt]1/3\\[2pt]-1/6
\end{pmatrix}.
\]

检查正交条件：

\[
A^{\trans}\hat r=
\begin{pmatrix}
\hat r_1+\hat r_2+\hat r_3\\
\hat r_1+2\hat r_2+3\hat r_3
\end{pmatrix}
=
\begin{pmatrix}
-1/6+1/3-1/6\\
-1/6+2(1/3)+3(-1/6)
\end{pmatrix}
=0.
\]

第一行说明残差之和为零——预测值和观测值的正负偏差相互抵消。第二行说明残差与时间坐标的加权和也为零——偏差中没有可以被“时间趋势”进一步解释的部分。这两条不是巧合：它们正是正规方程的两个分量，是每次正确的最小二乘拟合都自动满足的结构。

注意一件事：我们最小化的是预测值与观测值在 \(y\) 方向上的竖直距离，而不是每个点到直线的最短欧氏距离。前者对每个观测只惩罚竖直偏差；后者会同时惩罚横向偏差。这是两种不同的“近”，回到散点图时要分清。

\subsection{两幅几何图不要混为一谈}

最小二乘同一件事，可以画在两幅完全不同的几何图里。

第一幅在 \(\R^2\) 的散点图里：横轴是 \(t\)，纵轴是观测值 \(b_i\)。残差表现为数据点到拟合直线的竖直线段。最小二乘在最小化这些竖直距离的平方和。

第二幅在 \(\R^3\) 的向量空间里：观测被看成一个整体向量

\[
b=(b_1,b_2,b_3)^{\trans}\in\R^3,
\]

所有可能的预测构成 \(\operatorname{Col}(A)\) 中的一个平面。此时 \(\hat b\) 是 \(b\) 到该平面的正交投影，残差 \(\hat r\) 垂直于整个平面。

两幅图计算的是同一个量，但几何形态完全不同。散点图中的残差线段并不与拟合直线垂直；\(\R^3\) 中的残差向量则确实与列空间平面垂直。把这两者搞混，就会把“残差与 \(A\) 的列正交”误解成“散点图中的每个残差线段都与拟合直线成直角”——那是另一回事。

\subsection{从微积分也能走到同一个方程}

定义目标函数

\[
f(x)=\|Ax-b\|_2^2=(Ax-b)^{\trans}(Ax-b),
\]

考虑在 \(x\) 处加一个小扰动 \(h\)：

\[
\begin{aligned}
f(x+h)
&= \|Ax-b+Ah\|_2^2 \\
&= (Ax-b+Ah)^{\trans}(Ax-b+Ah) \\
&= (Ax-b)^{\trans}(Ax-b) + 2(Ax-b)^{\trans}Ah + (Ah)^{\trans}Ah \\
&= f(x) + 2(Ax-b)^{\trans}Ah + \|Ah\|_2^2.
\end{aligned}
\]

对 \(h\) 的一阶项是 \(2(Ax-b)^{\trans}Ah = 2(A^{\trans}(Ax-b))^{\trans}h\)，因此梯度为

\[
\nabla f(x)=2A^{\trans}(Ax-b).
\]

（矩阵微积分部分会系统处理这类求导；这里只需要看出一次项就够了。）令梯度为零，正是

\[
A^{\trans}A\hat x=A^{\trans}b.
\]

两条路线得出的方程完全一致。微积分说“这是一个凸二次函数，梯度为零处就是谷底”；投影说“残差与列空间正交处就是最近点”。前者便于推广到带约束和非线性的情况；后者直接揭示线性代数结构。知道两条路能通到同一个终点，比只会走其中一条更安心。

\subsection{什么是唯一的，什么不是}

无论如何选取参数，总存在唯一的预测向量 \(\hat b=A\hat x\)——有限维子空间中的正交投影总是唯一的。残差 \(\hat r=b-\hat b\) 也因此唯一。

但参数 \(\hat x\) 不一定唯一。如果 \(A\) 的两列线性相关（比如两列完全相同），模型实际上只依赖两个系数的某个组合，不同参数可能给出完全相同的预测向量。此时

\[
\hat x=(A^{\trans}A)^{-1}A^{\trans}b
\]

中的逆矩阵不存在，但最优拟合本身仍然在列空间里好好待着。

\(\hat x\) 唯一当且仅当 \(A\) 的列线性无关，这也等价于 \(A^{\trans}A\) 可逆。后面的 SVD 与伪逆部分会处理秩亏情形：从所有同样好的参数解中，选出范数最小的那一个。

\subsection{公式正确不等于计算方式理想}

当 \(A\) 满列秩时，正规方程的解形式上可以写成

\[
\hat x=(A^{\trans}A)^{-1}A^{\trans}b.
\]

但实际计算不应该先去求 \((A^{\trans}A)^{-1}\)。原因很具体：显式形成 \(A^{\trans}A\) 会把列之间的近线性相关成倍放大。在二范数意义下，\(A^{\trans}A\) 的条件数大约是 \(A\) 条件数的平方——条件数量化了“输入数据的微小相对扰动经过求解过程后，在输出中被放大多少倍的相对误差上限”（数值分析部分会给出严格定义）。

如果 \(A\) 的条件数已经不太好，平方之后可能直接让问题从“有点敏感”变成“不可靠”。

更稳妥的做法：

\begin{itemize}
\tightlist
\item
  满列秩的一般问题使用 QR 分解；
\item
  秩亏、近秩亏或需要最小范数解时使用 SVD；
\item
  只有确实需要对大量不同 \(b\) 反复做同一个投影时，才考虑显式使用投影矩阵。
\end{itemize}

下一节会展示 QR 怎么绕开 \(A^{\trans}A\) 的求逆，直接通过三角回代拿到参数。

\subsection{平方损失背后藏着一个假设}

平方损失对大残差的惩罚增长得很快——一个残差为 \(2\) 的观测对总损失的贡献是残差为 \(1\) 的观测的四倍。当误差大致对称、没有特别离谱的异常值时，这个特性很自然：让大偏差付出超额代价，会驱使我们优先缩小它们。

在线性高斯噪声模型中，最小二乘还恰好等价于最大似然估计（概率论与数理统计部分会回到这个精彩的对应）。但一个极端的离群点就可以把整条拟合线拽向自己，因为平方损失赋予了它不成比例的话语权。

如果换成绝对值损失，离群点的影响力会小得多；换成加权平方损失，可以手动调节不同观测的重要程度；换成 Huber 损失等稳健损失，则能在小残差区域保持平方行为、大残差区域转为线性。

用同一个三点数据快速对比一下。假如我们把第三个点从 \((3,2)\) 改成 \((3,10)\)——前两个点还在原来的位置，第三个点突然跳了上去。平方损失下的拟合直线会被这个离群点猛拽过去：正规方程变成

\[
3\hat C+6\hat D=13,\quad 6\hat C+14\hat D=32,
\]

解得 \(\hat D=3,\;\hat C=5/3\)——直线陡峭得多，几乎在迁就那个偏离的点。换成绝对值损失，第三个点的影响力不再按平方放大，直线会更靠近前两个点构成的趋势。损失函数的选择不是数学上的偏好，而是你在告诉模型：多大的偏差该承受多重的代价。每种损失函数都在定义一种不同的”近”，所以脱离了误差度量，不存在唯一的正确答案。

\subsection{一个极简的特例：均值也是最小二乘}

只用常数 \(C\) 拟合一组观测 \(b_1,\ldots,b_m\)。设计矩阵退化为一列全 \(1\)：

\[
A=\begin{pmatrix}1\\\vdots\\1\end{pmatrix}.
\]

正规方程变成

\[
A^{\trans}A\,\hat C = A^{\trans}b,
\]

即 \(m\hat C = \sum_{i=1}^m b_i\)，所以

\[
\hat C=\frac1m\sum_{i=1}^m b_i.
\]

样本均值就是在这个最简单的线性模型下的最小二乘估计。

翻回几何语言：从任意向量 \(b\in\R^m\) 中减去均值后，残差的各分量之和为零——这意味着残差与全 \(1\) 向量正交。减去均值就是先把 \(b\) 投影到全 \(1\) 方向上去掉的那一部分的补空间里。均值不是一个孤立的算术公式，它是投影到常数子空间后的坐标。

延伸阅读：MIT 18.06 Lecture 16：Projection Matrices and Least Squares。
